\documentclass{article}
\input{../../../../lib.sty}

\title{\texttt{fn make\_expr\_datetime\_component}}
\author{\MichaelShoemateAuthor}
\date{}

\begin{document}

\maketitle

\contrib
Proves soundness of \rustdoc{transformations/make\_stable\_expr/namespace\_dt/expr\_datetime\_component/fn}{make\_expr\_datetime\_component} in \asOfCommit{mod.rs}{f5bb719}.

\texttt{make\_expr\_datetime\_component} returns a Transformation that extracts a component from a temporal data type.

\section{Hoare Triple}
\subsection*{Precondition}
\subsubsection*{Compiler-verified}
\begin{itemize}
    \item Argument \texttt{input\_domain} of type \texttt{WildExprDomain}
    \item Argument \texttt{expr} of type \texttt{Expr}
    \item Argument \texttt{input\_metric} of type \texttt{M}
    \item Generic \texttt{M} implements \texttt{OuterMetric}
    \begin{itemize}
        \item \texttt{OuterMetric} defines an associated type \texttt{InnerMetric} that must implement \texttt{MicrodataMetric}
    \end{itemize}
    \item \texttt{(WildExprDomain, M)} and \texttt{(ExprDomain, M)} implement \texttt{MetricSpace}
    \item \texttt{M::Distance} implements \texttt{Clone}
    \item \texttt{Expr} implements \texttt{StableExpr<M, M>}
\end{itemize}

\subsubsection*{Caller-verified}
None

\subsection*{Pseudocode}
\lstinputlisting[language=Python,firstline=2,escapechar=`]{./pseudocode/make_expr_datetime_component.py}

\subsection*{Postcondition}
\begin{sloppypar}
\begin{theorem}
\validTransformation{\texttt{(input\_domain, input\_metric, expr, M)}}{\texttt{make\_expr\_datetime\_component}}
\end{theorem}
\end{sloppypar}
Starting from line \ref{match-strategy},
\texttt{expr} is matched as if it were a temporal expression, or otherwise rejects the expression.
All preconditions for \rustdoc{transformations/make\_stable\_expr/namespace\_dt/expr\_datetime\_component/fn}{match\_datetime\_component}
on line \ref{match-dt} are satisfied by the compiler,
therefore on success \texttt{to\_dtype} is the output data type of the recognized temporal component.

All preconditions for \texttt{make\_stable} on line \ref{t-prior} are compiler-verified,
therefore by the postcondition \texttt{t\_prior} is a valid transformation.
To prove that the output is a valid transformation,
we must first prove that the transformation on line \ref{t-component} is a valid transformation.

\begin{proof}[Proof of Data-Independent Errors]
    Line \ref{check-dtype} restricts the input to Time, Date, or nanosecond-resolution Datetime.
    The full signed 64-bit range of nanosecond timestamps is representable in Chrono;
    millisecond and microsecond timestamps can exceed that range and are rejected.
    Line \ref{check-component} rejects clock components on Date and calendar components on Time.
    These checks depend only on the public data type and selected operation.
    For the admitted combinations, extraction returns the component or null rather than raising data-dependent errors.
\end{proof}

\begin{proof}[Proof of Output Domain]
    The output domain is the same as the input domain,
    but the active series has a new data type and permits nulls.
    Line \ref{domain-dtype} updates the series domain so that non-null elements are of type \texttt{to\_dtype},
    which by the postcondition of \texttt{match\_datetime\_component} on line \ref{match-dt}
    represents the type of outputs when \texttt{temporal\_function} is applied to the data.

    Calendar extraction may introduce null values even on non-null input.
    Line \ref{nullable} therefore marks the output series as nullable.
    The name and context are preserved.
    The output domain now accounts for all transformations made to data in the input domain.
\end{proof}

\begin{proof}[Proof of Stability]
    Datetime component retrieval is a 1-stable row-by-row transformation,
    as the component retrieval is applied independently to each and every row.
    See \rustdoc{transformations/manipulation/fn}{make\_row\_by\_row} for a proof of the stability of row-by-row functions.
\end{proof}

Since it has been shown that both \texttt{t\_prior} and the component transformation are valid transformations,
then the preconditions for \texttt{make\_chain\_tt} are met (invoked via the right-shift operator shorthand),
and therefore the output is a valid transformation.

\end{document}
